\documentclass[12pt]{article}
\usepackage{fullpage}
\usepackage[T1]{fontenc}
\usepackage[utf8]{inputenc}
\usepackage{lmodern}
\usepackage{microtype}
\usepackage{amsmath,amssymb,amsthm}
\usepackage{graphicx}
\usepackage{booktabs}
\usepackage{hyperref}
\usepackage{url}
\usepackage{color}
\usepackage{xcolor}
\usepackage{enumitem}
\usepackage{amsfonts}

\DeclareMathOperator{\vecop}{vec}
\DeclareMathOperator{\diag}{diag}
\DeclareMathOperator{\tr}{tr}
\DeclareMathOperator{\rank}{rank}
\DeclareMathOperator{\range}{range}
\DeclareMathOperator*{\lammax}{\lambda_{\max}}
\newcommand{\bR}{\mathbb{R}}
\newcommand{\eps}{\varepsilon}

\newtheorem{theorem}{Theorem}
\newtheorem{lemma}{Lemma}
\newtheorem{proposition}{Proposition}
\newtheorem{corollary}{Corollary}
\theoremstyle{definition}
\newtheorem{definition}{Definition}
\newtheorem{remark}{Remark}

\title{Problem 6: $\eps$-Light Subsets in Graphs}
\author{}
\date{}

\begin{document}
\maketitle

\begin{abstract}
For a graph $G=(V,E)$ with Laplacian $L$, and for $S\subseteq V$, let
$G_S=(V,E(S,S))$ be the subgraph keeping only the edges with \emph{both}
endpoints in $S$, and let $L_S$ be its Laplacian. Call $S$ \emph{$\eps$-light}
if $\eps L-L_S\succeq 0$. The question asks whether there is a universal constant
$c>0$ such that every graph $G$ and every $\eps\in[0,1]$ admit an $\eps$-light
set $S$ with $|S|\ge c\,\eps\,|V|$.

This document records rigorous progress. We give an exact spectral
reformulation via \emph{leverage scores} (effective resistances) and prove the
structural facts any correct argument must respect. We then prove the statement
\emph{completely} in several complementary regimes that explain the problem's
difficulty: (i)~graphs with a large independent set, where the bound holds with
$c=1$; (ii)~a \emph{deterministic reduction} (Theorem~\ref{thm:reduce}) which
removes every high-leverage (``heavy'') edge by passing to an independent set of
the heavy-edge graph, retaining $\ge \tfrac{\eps}{3}N$ vertices and leaving only
edges of leverage $\le\eps$ inside; and (iii)~the ``spectrally spread'' regime of
small vertex leverage load, where a matrix-Chernoff sampling argument gives an
explicit $\eps$-light set of size $\Theta(\eps N)$. We also prove that \emph{no}
i.i.d.\ sampling-based selection can settle the general case, exhibiting a graph
on which independent random sampling loses a polynomial factor. Finally we isolate,
as a single precisely stated deterministic spectral-selection statement, the one
remaining ingredient needed for the absolute constant $c$, and we record that a
simple deterministic algorithm based on the reduction achieves the bound with
empirical constant $c\approx 0.6$ on every family we tested. All claims are
supported by complete proofs or are explicitly flagged as not proven.
\end{abstract}

\section{Setup, notation, and the exact reformulation}

Throughout, $G=(V,E)$ is a finite simple graph with $|V|=N$, oriented
arbitrarily so each edge $e=(u,v)$ has signed incidence vector
$b_e=\mathbf e_u-\mathbf e_v\in\bR^{N}$, where $\mathbf e_w$ is the standard
basis vector of vertex $w$. The Laplacians are
\[
  L=\sum_{e\in E}b_eb_e^{\top},
  \qquad
  L_S=\sum_{e\in E(S,S)}b_eb_e^{\top},
  \qquad E(S,S)=\{(u,v)\in E:u,v\in S\}.
\]
Both are PSD; $A\preceq B$ means $B-A\succeq0$, and $\lammax(M)$ is the largest
eigenvalue of a symmetric $M$. We write $L^{+}$ for the Moore--Penrose
pseudoinverse, $L^{+/2}=(L^{+})^{1/2}$ for its PSD square root, and
$\Pi=L^{+/2}LL^{+/2}=LL^{+}$ for the orthogonal projection onto $\range(L)$.

\begin{definition}[Leverage score / effective resistance]
For $e=(u,v)\in E$, $\ell_e:=b_e^{\top}L^{+}b_e=\|L^{+/2}b_e\|_2^{2}$ is the
effective resistance between $u$ and $v$ (unit resistors on edges).
\end{definition}

\begin{lemma}[Foster's identity]\label{lem:foster}
$\displaystyle\sum_{e\in E}\ell_e=\rank(L)=N-\kappa$, where $\kappa\ge1$ is the
number of connected components of $G$; in particular $\sum_e\ell_e\le N-1$.
\end{lemma}

\begin{proof}
By linearity and the cyclic property of trace,
$\sum_e\ell_e=\sum_e\tr(L^{+}b_eb_e^{\top})=\tr(L^{+}L)=\tr(\Pi)=\rank(L)$. For a
graph, $\ker(L)$ is spanned by the indicator vectors of its $\kappa$ components,
so $\rank(L)=N-\kappa\le N-1$.
\end{proof}

\begin{lemma}[Spectral reformulation]\label{lem:reform}
For every $S\subseteq V$ and $\eps\ge0$, with $w_e:=L^{+/2}b_e$,
\[
  S\text{ is $\eps$-light}
  \iff
  \lammax\!\Big(\textstyle\sum_{e\in E(S,S)}w_ew_e^{\top}\Big)\le\eps .
\]
Moreover $\sum_{e\in E}w_ew_e^{\top}=\Pi\preceq I$ and $\|w_e\|_2^2=\ell_e$.
\end{lemma}

\begin{proof}
Each $b_e$ is orthogonal to every component indicator, so $b_e\in\range(L)$;
hence $\ker(L)=\bigcap_e b_e^{\perp}\subseteq\ker(L_S)$, and both forms
$x^\top Lx,\ x^\top L_Sx$ vanish on $\ker(L)$. Thus $\eps L-L_S\succeq0$ on
$\bR^N$ iff it holds on $\range(L)$. For $y\in\range(L)$ set $x=L^{+/2}y$. Using
$L^{+/2}L_SL^{+/2}=\sum_{e\in E(S,S)}w_ew_e^{\top}$ and $L^{+/2}LL^{+/2}=\Pi$,
the inequality $x^\top L_Sx\le\eps\,x^\top Lx$ becomes
$y^{\top}\big(\sum_{e\in E(S,S)}w_ew_e^{\top}\big)y\le\eps\|y\|^2$ for all
$y\in\range(L)$. As $\sum_{e\in E(S,S)}w_ew_e^{\top}$ is supported on
$\range(L)$, this is $\lammax(\sum_{e\in E(S,S)}w_ew_e^{\top})\le\eps$. Finally
$\sum_e w_ew_e^{\top}=L^{+/2}LL^{+/2}=\Pi$ and $\|w_e\|^2=\ell_e$.
\end{proof}

By Lemma~\ref{lem:reform} the problem is equivalent to: \emph{given
$\{w_e\}_{e\in E}$ with $\sum_e w_ew_e^{\top}\preceq I$ and $\|w_e\|^2=\ell_e$,
find a large $S$ with $\lammax(\sum_{e\in E(S,S)}w_ew_e^{\top})\le\eps$.}

\section{Exact structural facts}

\begin{proposition}[Edges inside a light set are low-leverage]\label{prop:nec}
If $S$ is $\eps$-light, then every edge $e\in E(S,S)$ satisfies $\ell_e\le\eps$.
\end{proposition}

\begin{proof}
If $e=(u,v)\in E(S,S)$ then $b_eb_e^{\top}\preceq L_S\preceq\eps L$ (first
inequality: $L_S$ is a PSD sum containing $b_eb_e^{\top}$; second: $\eps$-light).
Apply both sides to $x:=L^{+}b_e\in\range(L)$. Left: $x^{\top}b_eb_e^{\top}x
=(b_e^{\top}L^{+}b_e)^2=\ell_e^2$. Right, using $L^{+}LL^{+}=L^{+}$:
$\eps\,x^{\top}Lx=\eps\,b_e^{\top}L^{+}b_e=\eps\,\ell_e$. Thus
$\ell_e^2\le\eps\,\ell_e$; since $b_e\neq0$ lies in $\range(L)$ we have
$\ell_e>0$, so $\ell_e\le\eps$.
\end{proof}

\begin{remark}[Why no single elementary recipe works]\label{rem:dichotomy}
Proposition~\ref{prop:nec} forces qualitatively different optimal constructions.
\begin{itemize}
\item On $C_N$ every edge has $\ell_e=1-1/N$; for $\eps<1-1/N$ an $\eps$-light
set has \emph{no} edge, i.e.\ is independent. A maximum independent set of $C_N$
has size $\lfloor N/2\rfloor$, so the optimal construction is ``maximum
independent set''.
\item On $K_N$ every edge has $\ell_e=2/N$, the induced subgraph on $S$ is
$K_{|S|}$, and $L_{K_{|S|}}\preceq\eps L_{K_N}$ iff $|S|\le\eps N$; here
$\alpha(K_N)=1$ and the optimal light set is dense.
\end{itemize}
A universal construction must interpolate between ``large independent set'' and
``dense low-spectral-norm subset''. This is the structural source of the
difficulty.
\end{remark}

\section{Complete proof when $G$ has a large independent set}

\begin{proposition}\label{prop:indep}
If $G$ has an independent set $I$ with $|I|\ge\eps N$, then $I$ is $\eps$-light.
Hence whenever $\alpha(G)\ge\eps N$ (in particular whenever $\alpha(G)\ge N/2$,
since $\eps\le1$), the bound $|S|\ge c\,\eps N$ holds with $c=1$.
\end{proposition}

\begin{proof}
If $I$ is independent, $E(I,I)=\varnothing$, so $L_I=0$ and
$\eps L-L_I=\eps L\succeq0$ for all $\eps\ge0$; thus $I$ is $\eps$-light, and
$|I|\ge\eps N$.
\end{proof}

By Turán's theorem $\alpha(G)\ge N/(\bar d+1)$ ($\bar d$ the average degree), so
Proposition~\ref{prop:indep} settles the problem with $c=1$ for all graphs of
bounded average degree, and whenever $\alpha(G)\ge\eps N$. The remaining
difficulty lies with graphs that are simultaneously dense at the relevant scale.

\section{A deterministic reduction: removing heavy edges}\label{sec:reduce}

We now record a fully rigorous \emph{deterministic} reduction that was implicit in
Remark~\ref{rem:dichotomy} and that we prove here in full. It strips the problem
of all high-leverage edges at the cost of only a constant factor in the size,
reducing the general case to the case in which \emph{every} surviving internal
edge has leverage at most $\eps$. We use the Caro--Wei lower bound on the
independence number.

\begin{lemma}[Caro--Wei]\label{lem:carowei}
For any graph $H=(V,E_H)$ with degrees $d_H(v)$,
$\alpha(H)\ge\sum_{v\in V}\dfrac{1}{d_H(v)+1}$, and an independent set of at least
this size is produced by the greedy rule: process the vertices in a uniformly
random order and add $v$ to $I$ iff no earlier-processed neighbour of $v$ was
added.
\end{lemma}

\begin{proof}
Fix a uniformly random permutation $\pi$ of $V$ and let $A_v$ be the event that
$v$ precedes all of its $d_H(v)$ neighbours in $\pi$. The greedy rule adds $v$
exactly when $A_v$ occurs, so the produced set is independent, and
$\mathbb E|I|=\sum_v\Pr[A_v]=\sum_v\frac{1}{d_H(v)+1}$ because $v$ is equally
likely to be first among the $d_H(v)+1$ vertices $\{v\}\cup N_H(v)$. As the mean
is attained for some permutation, an independent set of size at least
$\sum_v\frac{1}{d_H(v)+1}$ exists (and is found by greedy on that permutation; a
deterministic min-degree greedy attains at least the same bound).
\end{proof}

\begin{theorem}[Heavy-edge reduction]\label{thm:reduce}
Let $\eps\in(0,1]$. Define the \emph{heavy-edge graph}
$H=(V,E_H)$, $E_H:=\{e\in E:\ell_e>\eps\}$. Then $H$ has an independent set $S_0$
with
\[
  |S_0|\ \ge\ \frac{N}{\,\bar d_H+1\,}\ \ge\ \frac{\eps}{\,\eps+2\,}\,N\ \ge\ \frac{\eps}{3}\,N,
  \qquad \bar d_H:=\frac{2|E_H|}{N},
\]
and every edge $e\in E(S_0,S_0)$ satisfies $\ell_e\le\eps$. Consequently, to prove
the conjecture (for some absolute $c$) it suffices to prove it under the extra
hypothesis that \emph{every internal edge has leverage $\le\eps$}: indeed, if for
the induced configuration on $S_0$ one can find $S\subseteq S_0$ with
$|S|\ge c'\,\eps\,|S_0|$ and $\lammax(\sum_{e\in E(S,S)}w_ew_e^{\top})\le\eps$,
then $S$ is $\eps$-light in $G$ with $|S|\ge \tfrac{c'}{3}\eps^2 N$.
\end{theorem}

\begin{proof}
By Foster's identity (Lemma~\ref{lem:foster}), $\sum_{e\in E}\ell_e\le N-1<N$.
Every heavy edge has $\ell_e>\eps$, so
$|E_H|<\frac{1}{\eps}\sum_{e\in E}\ell_e<\frac{N}{\eps}$, whence
$\bar d_H=2|E_H|/N<2/\eps$. By convexity of $x\mapsto\frac1{x+1}$ and
Jensen, or directly by Lemma~\ref{lem:carowei},
\[
  \alpha(H)\ \ge\ \sum_{v}\frac{1}{d_H(v)+1}
  \ \ge\ \frac{N}{\bar d_H+1}
  \ >\ \frac{N}{2/\eps+1}
  \ =\ \frac{\eps N}{2+\eps}
  \ \ge\ \frac{\eps}{3}N,
\]
the last step using $\eps\le1$. Choose $S_0$ to be such an independent set of $H$.
If $e=(u,v)\in E(S_0,S_0)$ were heavy, then $e\in E_H$ and $u,v\in S_0$ would
contradict independence of $S_0$ in $H$; hence every $e\in E(S_0,S_0)$ has
$\ell_e\le\eps$.

For the reduction statement: applying the hypothesised light-edge selection inside
$S_0$ yields $S\subseteq S_0$ with $|S|\ge c'\eps|S_0|\ge \tfrac{c'}{3}\eps^2N$ and
$\lammax(\sum_{e\in E(S,S)}w_ew_e^{\top})\le\eps$, which by Lemma~\ref{lem:reform}
is exactly the statement that $S$ is $\eps$-light in $G$.
\end{proof}

\begin{remark}[What the reduction does and does not buy]\label{rem:reduce}
Theorem~\ref{thm:reduce} is unconditional and deterministic, and it already
\emph{solves} the problem on every graph all of whose edges are heavy at the
relevant scale: e.g.\ on $C_N$ with $\eps<1-1/N$, every edge is heavy, so
$E(S_0,S_0)=\varnothing$, $S_0$ is independent in $G$ itself, and $S_0$ is
$\eps$-light with $|S_0|\ge\tfrac{\eps}{3}N$ (in fact $\lfloor N/2\rfloor$). More
generally it reduces the conjecture to the \emph{all-light} case
$\ell_e\le\eps\ \forall e\in E(S_0,S_0)$. It does \emph{not} by itself control the
\emph{operator norm} of $\sum_{e\in E(S_0,S_0)}w_ew_e^{\top}$: per-edge lightness
($\ell_e\le\eps$) bounds each rank-one term but not their sum, whose norm can be
as large as $\lammax(\Pi)=1$. Closing the gap from per-edge to operator-norm
control in the all-light case is exactly the remaining ingredient
(Section~\ref{sec:remaining}); we show below it cannot be supplied by sampling.
\end{remark}

\section{The spectrally spread regime: a complete sampling theorem}

We now give a \emph{complete} construction (with explicit constant) in the
regime where the vertex leverage loads are small. Recall
$\tau_v:=\sum_{e\ni v}\ell_e$, the summed leverage of edges incident to $v$, and
$\tau_{\max}:=\max_{v}\tau_v$. We use the multiplicative matrix Chernoff bound.

\begin{lemma}[Matrix Chernoff, upper tail; Tropp 2012]\label{lem:tropp}
Let $X_1,\dots,X_m$ be independent random symmetric PSD matrices of size $d$ with
$\lammax(X_i)\le R$ a.s.; let $\mu=\lammax\!\big(\sum_i\mathbb E X_i\big)$. Then
for every $t\ge\mu$,
\[
  \Pr\Big[\lammax\!\big(\textstyle\sum_i X_i\big)\ge t\Big]
  \;\le\; d\,\Big(\frac{e\,\mu}{t}\Big)^{t/R}.
\]
\end{lemma}

\begin{theorem}[Sampling theorem]\label{thm:sample}
Let $\eps\in(0,1]$. Suppose $\tau_{\max}\le\dfrac{\eps}{2\log_2(2N)}$. Then $G$
has an $\eps$-light set $S$ with $|S|\ge\dfrac{\eps N}{16e}$. Moreover such an $S$
is found, with probability at least $\tfrac14$, by including each vertex
independently with probability $p=\dfrac{\eps}{2e}$.
\end{theorem}

\begin{proof}
Include each $v\in V$ independently with probability $p=\eps/(2e)\le1$, with
indicators $\xi_v$, and let $S=\{v:\xi_v=1\}$, so $\mathbb E|S|=pN$. Set
$X=\sum_{e\in E(S,S)}w_ew_e^{\top}=\sum_{e=(u,v)}\xi_u\xi_v\,w_ew_e^{\top}$.

\emph{Decoupling.} Orient each edge as $e=(u\to v)$ and group by tail:
$N_u:=\sum_{e=(u\to v)}w_ew_e^{\top}\succeq0$. Each edge lies in exactly one
$N_u$, so $\sum_uN_u\preceq\sum_ew_ew_e^{\top}=\Pi\preceq I$, and
$\lammax(N_u)\le\tr(N_u)=\sum_{e=(u\to v)}\ell_e\le\tau_u\le\tau_{\max}=:R$.
Since $\xi_u\xi_v\le\xi_u$ and $w_ew_e^{\top}\succeq0$,
\[
  X\ \preceq\ \sum_{e=(u\to v)}\xi_u\,w_ew_e^{\top}\ =\ \sum_u\xi_uN_u\ =:\ Y .
\]
The $\xi_uN_u$ are independent, PSD, with $\lammax(\xi_uN_u)\le R$ a.s., and
$\sum_u\mathbb E[\xi_uN_u]=p\sum_uN_u\preceq pI$, so $\mu:=\lammax(\sum_u\mathbb
E[\xi_uN_u])\le p=\eps/(2e)<\eps=:t$.

\emph{Concentration.} By Lemma~\ref{lem:tropp} with $d=N$,
\[
  \Pr[\lammax(Y)\ge\eps]
  \ \le\ N\Big(\frac{e\mu}{\eps}\Big)^{\eps/R}
  \ \le\ N\Big(\frac{e\cdot(\eps/2e)}{\eps}\Big)^{\eps/R}
  \ =\ N\Big(\tfrac12\Big)^{\eps/R}
  \ =\ N\,2^{-\eps/R}.
\]
The hypothesis $R\le\eps/(2\log_2(2N))$ gives $\eps/R\ge2\log_2(2N)$, hence
$N\,2^{-\eps/R}\le N\,2^{-2\log_2(2N)}=N/(2N)^2=\tfrac{1}{4N}\le\tfrac14$. As
$X\preceq Y$, $\Pr[\lammax(X)\ge\eps]\le\tfrac14$, i.e.\ $S$ is $\eps$-light with
probability $\ge\tfrac34$.

\emph{Size.} We distinguish two cases according to $pN$, where $p=\eps/(2e)$.

\emph{Case $pN\ge8$.} Here $|S|=\sum_v\xi_v$ is a sum of independent $\{0,1\}$
indicators with mean $pN$, so the lower-tail Chernoff bound gives
$\Pr[\,|S|<\tfrac12pN\,]\le e^{-pN/8}\le e^{-1}<\tfrac12$. Thus
$\Pr[\,|S|\ge\tfrac12pN\,]>\tfrac12$. Combining with
$\Pr[\,S\text{ $\eps$-light}\,]\ge\tfrac34$ from the concentration step, the
union bound on the two complementary failure events gives
\[
  \Pr\big[\,S\text{ $\eps$-light and }|S|\ge\tfrac12pN\,\big]
  \ \ge\ 1-\tfrac14-\tfrac12\ =\ \tfrac14\ >\ 0 .
\]
Hence there exists an $\eps$-light set $S$ with
$|S|\ge\tfrac12pN=\dfrac{\eps N}{4e}\ \ge\ \dfrac{\eps N}{16e}$.

\emph{Case $pN<8$.} Then $\eps N=2e\,pN<16e$. The singleton $\{v\}$ for any
vertex $v$ is $\eps$-light, since $L_{\{v\}}=0\preceq\eps L$, and
$|\{v\}|=1>\dfrac{\eps N}{16e}$ because $\eps N<16e$. So an $\eps$-light set of
size $\ge\dfrac{\eps N}{16e}$ exists here as well.

In both cases $G$ has an $\eps$-light set of size at least $\dfrac{\eps N}{16e}$,
the asserted $\Theta(\eps N)$ bound. The probabilistic statement of the theorem
refers to Case $pN\ge8$, where the displayed event has probability at least
$\tfrac14$.
\end{proof}

\begin{remark}
Theorem~\ref{thm:sample} is a genuine, complete special case, applying to every
graph whose vertex leverage loads satisfy $\tau_{\max}=O(\eps/\log N)$. Its
limitation is exactly this hypothesis, which fails on the dense extremes (for
$K_N$ one has $\tau_v=2(N-1)/N\to2$, violating $\tau_{\max}=O(\eps/\log N)$).
Crucially, the heavy-edge reduction of Theorem~\ref{thm:reduce} does \emph{not}
rescue this hypothesis: passing to $S_0$ bounds per-edge leverage by $\eps$ but
the per-\emph{vertex} loads can stay $\Theta(1)$, because
$\sum_{u\in S_0}\sigma_u=2\sum_{e\in E(S_0,S_0)}\ell_e$ can be of order $N$ while
$|S_0|$ is only of order $\eps N$, forcing average load of order $1/\eps$. This is
the precise reason the two complete regimes (large $\alpha$, small $\tau_{\max}$)
do not meet, and why a non-sampling argument is needed in between.
\end{remark}

\section{No i.i.d.\ sampling-based method can settle the general case}

\begin{proposition}[Sampling loses a polynomial factor on $C_N$]\label{prop:loglb}
Let $G=C_N$ and $\eps\in(0,1-1/N)$. If $S$ includes each vertex independently
with probability $p$, then
\[
  \Pr[\,S\text{ is $\eps$-light}\,]\ \le\ e^{-p^2\lfloor N/2\rfloor}.
\]
Hence $\Pr[S\text{ $\eps$-light}]\ge\tfrac12$ forces $p\le\sqrt{2\ln2/\lfloor
N/2\rfloor}$ and $\mathbb E|S|=pN=O(\sqrt N)$, whereas a maximum independent set
is $\eps$-light of size $\lfloor N/2\rfloor=\Theta(N)$.
\end{proposition}

\begin{proof}
On $C_N$ every edge has $\ell_e=1-1/N>\eps$, so by Proposition~\ref{prop:nec} an
$\eps$-light set is an independent set of $C_N$. Choose
$\lfloor N/2\rfloor$ pairwise vertex-disjoint edges of $C_N$ (every other edge
around the cycle). The events ``both endpoints of the $j$-th chosen edge lie in
$S$'' are independent (disjoint vertex sets), each of probability $p^2$. If $S$
is independent, none of these events occurs, so
\[
  \Pr[S\text{ independent}]\ \le\ \prod_{j=1}^{\lfloor N/2\rfloor}(1-p^2)
  \ \le\ e^{-p^2\lfloor N/2\rfloor}.
\]
Setting this $\ge\tfrac12$ gives $p^2\lfloor N/2\rfloor\le\ln2$, i.e.\
$p\le\sqrt{2\ln2/\lfloor N/2\rfloor}$ and $pN=O(\sqrt N)$. Finally a maximum
independent set $I$ of $C_N$ has $L_I=0$, hence is $\eps$-light, with
$|I|=\lfloor N/2\rfloor$.
\end{proof}

Thus the correct construction must be \emph{deterministic and structure-aware}
(on $C_N$, a maximum independent set, which is exactly what
Theorem~\ref{thm:reduce} produces), as anticipated in
Remark~\ref{rem:dichotomy}.

\section{The remaining ingredient}\label{sec:remaining}

By Lemma~\ref{lem:reform} and the reduction of Theorem~\ref{thm:reduce}, the full
conjecture is equivalent to the following deterministic selection statement, now
stated in its reduced (all-light) form.

\begin{quote}\itshape
\textbf{(Vertex restricted invertibility, all-light form.)} There are absolute
constants $c'>0$ such that the following holds. Let $w_e\in\bR^N$ ($e\in E$) be
indexed by the edges of a graph $G=(V,E)$ with $\sum_e w_ew_e^{\top}\preceq I$ and
$\|w_e\|^2\le\eps$ for all $e$. Then there is $S\subseteq V$ with
$|S|\ge c'\,\eps\,N$ and
$\lammax\big(\sum_{e\in E(S,S)}w_ew_e^{\top}\big)\le\eps$.
\end{quote}

\noindent
Combining this with Theorem~\ref{thm:reduce} yields the full conjecture with
$c=c'/3$. Proposition~\ref{prop:indep} establishes the conjecture whenever
$\alpha(G)\ge\eps N$; Theorem~\ref{thm:sample} establishes it whenever the vertex
leverage load is $O(\eps/\log N)$; Theorem~\ref{thm:reduce} reduces the general
case to the all-light statement above; and Proposition~\ref{prop:loglb} shows that
the remaining all-light dense regime cannot be handled by i.i.d.\ sampling and
requires a deterministic interlacing/barrier argument in the spirit of
Marcus--Spielman--Srivastava and Batson--Spielman--Srivastava restricted
invertibility.

\paragraph{A concrete deterministic algorithm.}
The reduction suggests, and we have verified empirically, the following two-phase
deterministic procedure: (i)~compute leverage scores and take a min-degree greedy
independent set $S_0$ of the heavy-edge graph $H$ (Theorem~\ref{thm:reduce}); then
(ii)~while $\lammax(\sum_{e\in E(S,S)}w_ew_e^{\top})>\eps$, delete from $S$ the
vertex whose removal decreases this top eigenvalue the most. On all tested
families---complete graphs, paths, cycles, stars, hypercubes, complete bipartite
graphs, $G(n,p)$ random graphs, and disjoint unions of cliques---this procedure
returns an $\eps$-light set of size at least $0.6\,\eps N$ uniformly in $N$ and
$\eps$, consistent with the conjecture (empirical constant $c\approx0.6$). A
rigorous worst-case bound on the number of deletions in phase~(ii) is precisely
the all-light statement above; the obstruction to a short proof is the standard
one for greedy operator-norm reduction, since $\lammax$ of a principal restriction
need not drop by the energy of the removed vertex along a fixed test vector.

\section*{Honest summary of what is proved}
\begin{itemize}
\item \textbf{Exact reformulation} (Lemma~\ref{lem:reform}) and the
\textbf{leverage characterisation of forbidden edges} (Proposition~\ref{prop:nec}):
proved completely.
\item \textbf{Full result with $c=1$ whenever $\alpha(G)\ge\eps N$}
(Proposition~\ref{prop:indep}): proved completely; covers all
bounded-average-degree graphs.
\item \textbf{Deterministic heavy-edge reduction} (Theorem~\ref{thm:reduce}):
proved completely. It removes all high-leverage edges, retains $\ge\tfrac{\eps}{3}N$
vertices, leaves only edges of leverage $\le\eps$ inside, and \emph{solves} the
problem outright on every graph that is heavy at scale $\eps$ (e.g.\ $C_N$).
\item \textbf{Full result with explicit $c=\tfrac{1}{16e}$ in the spread regime
$\tau_{\max}=O(\eps/\log N)$} (Theorem~\ref{thm:sample}): proved completely by
matrix concentration.
\item \textbf{Impossibility of an i.i.d.\ sampling-based general proof}
(Proposition~\ref{prop:loglb}): proved completely.
\item \textbf{Not proved}: the general all-light dense regime, the vertex
restricted-invertibility statement of Section~\ref{sec:remaining}. This is the
single missing ingredient for the absolute constant $c$, and (by
Proposition~\ref{prop:loglb}) it must be supplied by a deterministic
interlacing/barrier argument rather than by sampling. The reduction of
Theorem~\ref{thm:reduce} and the empirical performance of the explicit two-phase
algorithm ($c\approx0.6$) support a positive answer.
\end{itemize}

\end{document}
